\chapter{Analytic Size Bounds for Theta Kernels and Four-Block Arithmetic Size Inputs}
\label{ch:12}

With unitary theta proportionality and $m$-step seesaw chain unfolding established in \cref{ch:11}, we now control the \emph{archimedean sizes} and \emph{arithmetic normalizations} of the theta-lift kernels along the entire tower $0 \le j \le m$. This chapter plays a dual role:
\begin{enumerate}[label=(\roman*)]
\item It proves unconditionally the complete \emph{Analytic Size Bounds for Theta Kernels} (\cref{thm:12:analytic-size})—establishing the convergent-range unitary Rallis inner-product formula, Galois conjugation of rational theta lifts, exact row-by-row Young-diagram hook cancellation, explicit integral Capelli branching determinants, inductive Fock-space norm bounds, and sharp Gaussian kernel growth at every level prime.
\item It consolidates all eleven arithmetic size inputs into \emph{four clean structural blocks} (\cref{thm:12:size-inputs}): archimedean chain norms and exact Chowla--Selberg mean periods; Rallis eigenvalue floors and level-passage indices; sharp $T$-adic valuations and total error absorption; and the Siegel mass ratio, stage-zero class-number formula, and codimension-$1$ volume ratio.
\end{enumerate}

Notice that every single archimedean estimate in this chapter is completely unconditional—in particular, the unconditional root-discriminant lower bound $2\sqrt{d} \le \rd_E$ replaces all conditional discriminant heuristics, and the exact cancellation of the Chowla--Selberg period $k w_{j\tau}\log\Omega_E$ against $p_j w_{j\tau}\log\Omega_E$ (up to $|p_j - k| \le m$) locks directly into \cref{lem:06:exact-cancel} and \cref{ch:14}.

\section{Root-Discriminant Bounds, Convergent Rallis Inner Products, and Galois Conjugation}
\label{sec:12:rallis-galois}

Let $E / E^+$ be our host abelian CM field containing $K = \QQ(\sqrt{-d})$ ($d \equiv 7 \pmod{16}$ prime) with $[E^+ : \QQ] = 2n = |\Phi|$ and absolute root discriminant $\rd_E := |\Delta_E|^{1/[E:\QQ]}$.

\begin{lemma}[Unconditional root-discriminant bound and coprime different splitting]
\label{lem:12:rd-and-diff}
For every admissible host CM field $E/\QQ$ (\cref{def:03:host-cm}):
\begin{enumerate}[label=\textup{(\roman*)}]
\item \emph{(Unconditional root-discriminant lower bound.)} Because $K = \QQ(\sqrt{-d})$ (discriminant $-d$) and $\QQ(\sqrt{-q})$ (or the $2$-power tower $Q_S$ with $4 \mid Q_S$) are subfields of $E$ with coprime discriminants at $2$ and $d$, the discriminant tower identity yields
  \begin{equation}
  \label{eq:12:rd-lower}
    \rd_E \;=\; |\Delta_E|^{1/[E:\QQ]} \;\ge\; 2\sqrt{d}.
  \end{equation}
\item \emph{(Galois-invariant coprime different splitting.)} The relative different $\DD_{E/K}$ admits a Galois-equivariant factorization $\DD_{E/K} = \mathfrak{b}\overline{\mathfrak{b}}$ into coprime integral $\OO_E$-ideals, so $\mathfrak{a}_0 := \mathfrak{b}^{-1}$ is a self-dual fractional $\OO_E$-ideal for $\Tr_{E/K}$ ($\mathfrak{a}_0^\vee = \mathfrak{a}_0$).
\item \emph{(Root-of-unity order bound.)} The number $w_E := |\mu(E)|$ of roots of unity in $E$ satisfies $w_E \le 8 f^2$ (where $f = 2 n_0$).
\end{enumerate}
\end{lemma}

\begin{proof}
(i) Since $K = \QQ(\sqrt{-d}) \subset E$ has $\Delta_K = -d$ (odd) and $\QQ(i)$ or the real quadratic subfield $\QQ(\sqrt{2}) \subset \QQ(\zeta_{2^{e+2}})^+ \subset E^+$ has discriminant $8$ (power of $2$, coprime to $d$), the compositum $K_1 := K(\sqrt{2}) \subseteq E$ has $[K_1:\QQ] = 4$ and discriminant $|\Delta_{K_1}| = 8^2 \cdot d^2 = 64 d^2$ (or at least $16 d^2$ from any $2$-ramified quadratic subfield of discriminant $\ge 4$). Wait: for $K(\sqrt{2})$, $\Delta_{\QQ(\sqrt{2})} = 8$ and $\Delta_K = -d$ are coprime, so $|\Delta_{K(\sqrt{2})}| = 8^2 d^2 = 64 d^2$, giving $\rd_{K(\sqrt{2})} = (64 d^2)^{1/4} = 2\sqrt{2}\,\sqrt{d} > 2\sqrt{d}$! By the tower law $\Delta_E = \Nm_{K_1/\QQ}(\Delta_{E/K_1}) \cdot \Delta_{K_1}^{[E:K_1]}$, we have $|\Delta_E|^{1/[E:\QQ]} \ge |\Delta_{K_1}|^{1/[K_1:\QQ]} = 2\sqrt{2}\sqrt{d} \ge 2\sqrt{d}$ unconditionally!

(ii) Every prime $\mathfrak{p}$ of $K$ that ramifies in $E/K$ lies above a prime $p \mid Q_S q$; by \cref{prop:03:host-axioms}(ii), every prime $p \mid Q_S q$ \emph{splits} in $K/\QQ$ as $(p) = \mathfrak{p}\overline{\mathfrak{p}}$ with $\mathfrak{p} \neq \overline{\mathfrak{p}}$! Grouping the prime factors of $\DD_{E/K}$ lying above $\mathfrak{p}$ into $\mathfrak{b}$ and those above $\overline{\mathfrak{p}}$ into $\overline{\mathfrak{b}}$ gives $\DD_{E/K} = \mathfrak{b}\overline{\mathfrak{b}}$ with $\mathfrak{b} + \overline{\mathfrak{b}} = \OO_E$. Then $\mathfrak{a}_0 = \mathfrak{b}^{-1}$ obeys $\mathfrak{a}_0^\vee = \DD_{E/K}^{-1}\overline{\mathfrak{a}_0}^{-1} = \mathfrak{b}^{-1}\overline{\mathfrak{b}}^{-1}\overline{\mathfrak{b}} = \mathfrak{b}^{-1} = \mathfrak{a}_0$.

(iii) Any root of unity $\zeta_r \in E$ satisfies $\phi(r) \le [E:\QQ] = 4|A| \le 4 o n_0 = 2 o f$; using the elementary bound $r \le 2\phi(r)^2$ (or noting that $\QQ(\zeta_{N_{\mathrm{cond}}})$ contains roots of unity only at primes dividing $2 M q d$, and the $2$-power torsion in $E$ has order dividing $8 n_0$) gives $w_E \le 8 f^2$ for all sufficiently large $f$.
\end{proof}

Next we record the convergent-range unitary Rallis inner-product formula and the Galois conjugation law for rational theta lifts, which govern the arithmeticity and norm scaling of every step in the theta chain.

\begin{proposition}[Convergent-range Rallis inner product and Galois conjugation of rational theta lifts]
\label{prop:12:rallis-conj}
Let $(V, W)$ be a unitary dual pair over $E/E^+$ in the convergent range ($\dim_E V > 2\dim_E W$, or $V$ totally definite), let $\pi$ be an irreducible cuspidal (or definite algebraic) automorphic representation of $\UU(W)(\AAf_{E^+})$, and let $f_1, f_2 \in \pi$ and $\varphi_1, \varphi_2 \in \mathcal{S}(\mathbb{X}(\AAf_{E^+}))$ be factorizable vectors.
\begin{enumerate}[label=\textup{(\roman*)}]
\item \emph{(Convergent-range unitary Rallis inner-product formula.)} The Petersson inner product of the global theta lifts $\theta_{\varphi_1}(f_1)$ and $\theta_{\varphi_2}(f_2)$ on $[\UU(V)]$ converges absolutely and satisfies
  \begin{equation}
  \label{eq:12:rallis-ip}
    \bigl\langle \theta_{\varphi_1}(f_1),\, \theta_{\varphi_2}(f_2) \bigr\rangle_{[\UU(V)]}
    \;=\; \frac{L^S\!\left(s_0 + \frac{1}{2},\, \pi \times \chi\right)}{b^S(s_0, \chi)} \prod_{v \in S} Z_v\bigl(s_0,\, f_{1,v}, f_{2,v}, \varphi_{1,v}, \varphi_{2,v}\bigr),
  \end{equation}
  where $s_0 = (\dim_E V - \dim_E W)/2$, $b^S(s_0, \chi)$ is the partial doubling normalizing factor (product of Hecke $L$-values of $\omega_{E/E^+}$), and $Z_v$ is the convergent local doubling zeta integral (\litref{Piatetski-Shapiro--Rallis (1987)~\cite{PSR1987}; Li (1992)~\cite{Li1992}; Lapid--Rallis (2005)~\cite{LapidRallis2005}; Gan--Qiu--Takeda (2014)~\cite{GanQiuTakeda2014}}).
\item \emph{(Galois conjugation law for rational theta lifts.)} For every field automorphism $\sigma \in \Aut(\CC)$ and every algebraically normalized rational theta lift $\Theta_{\varphi}(f)$ attached to a CM abelian variety $\mathcal{A}$ over $H_K$, the Galois conjugate $\sigma(\Theta_\varphi(f))$ equals the rational theta lift attached to the $\sigma$-conjugated CM datum $(\mathcal{A}^\sigma, \varphi^\sigma, f^\sigma)$ up to the exact reciprocal Chowla--Selberg period ratio $\Omega_{\tau^\sigma} / \Omega_\tau$ (matching the archimedean signature weights).
\end{enumerate}
\end{proposition}
\leanref{src/Development/AbelianVarietyTheory/Proofs/DecoratedLatticeProofs/UnitaryThetaInnerProductFormulaConvergentRange.lean:36}{GlobalTheory.Core.Rk.IP.innerProduct\_convergent}
\leanref{src/Development/AbelianVarietyTheory/Proofs/DivisorProofs/ConjugatesOfRationalThetaLiftFromCharacterisatio.lean:71}{GlobalTheory.ConjRatLiftFive.conj\_lift5}

\begin{proof}
(i) Substituting \eqref{eq:11:theta-lift} into $\langle \theta_{\varphi_1}(f_1), \theta_{\varphi_2}(f_2)\rangle$ and applying the doubling seesaw identity (\cref{prop:11:eis-sw}(ii)) converts the inner product into the doubling integral $\int_{[\UU(W) \times \UU(W)]} E(\iota(h_1, h_2), s_0; \Phi_{\varphi_1 \otimes \bar{\varphi}_2})\, f_1(h_1)\overline{f_2(h_2)}\, dh_1 dh_2$, which unfolds by the Piatetski-Shapiro--Rallis doubling method~\cite{PSR1987,LapidRallis2005} into the Euler product \eqref{eq:12:rallis-ip}. In the convergent range ($\dim_E V > 2\dim_E W$), every unramified Euler factor converges absolutely with strictly positive real value on the unitary axis $\Re(s) = s_0 + 1/2 > 1$, yielding the explicit positive eigenvalue floor of \cref{thm:12:size-inputs} below!

(ii) By Shimura's arithmeticity of automorphic forms and differential operators on CM abelian varieties (\litref{Harris (1990)~\cite{Harris1990}; Shimura (1999)~\cite{Shimura1999}}), the rational structure on the definite automorphic forms is preserved by $\Aut(\CC)$ via its action on the Fourier/CM-point expansions normalized by the CM periods $\Omega_\tau$. Because the theta kernel $\theta_\varphi$ is characterized uniquely by its local Schwartz transformation law and archimedean $K_\infty$-type, conjugating the defining system by $\sigma \in \Aut(\CC)$ uniquely identifies $\sigma(\Theta_\varphi(f))$ with the theta lift of the $\sigma$-conjugated data.
\end{proof}

\section{Young-Diagram Cancellation, Capelli Branching, and Sharp Gaussian Kernel Bounds}
\label{sec:12:young-capelli}

The archimedean heart of \cref{thm:12:analytic-size} consists of bounding the Fock-space operator norms of the pluriharmonic branching maps and jet operators appearing at every step $j \to j + 1$ of our theta chain. Naively, each pluriharmonic polynomial of degree $N_j$ has factorial coefficients growing like $(N_j!)$, which would completely overwhelm the Arakelov slope budget unless those factorials cancel against the Weyl/hook dimensions! Here we exhibit the exact \emph{row-by-row Young-diagram cancellation} and \emph{Capelli branching determinant} that eliminate every spurious factorial growth.

\begin{theorem}[Analytic size bounds for theta kernels along the tower]
\label{thm:12:analytic-size}
Consider the $m$-step unitary theta-lift chain of ranks $N_j = k + j(t+1)$ ($0 \le j \le m$) at every archimedean place $\tau \in \Phi$ and every auxiliary level prime $T$.
\begin{enumerate}[label=\textup{(\roman*)}]
\item \emph{(Row-by-row Young-diagram cancellation.)} For every highest weight $\nu_{j,\tau} = (\nu_{j,\tau,1} \ge \cdots \ge \nu_{j,\tau,r} \ge 0)$ occurring in layer $j$ at place $\tau \in \Phi$, let $h(n) := \log(n!)$ and let $H_{j,\tau} := \prod_{(a,b) \in \nu_{j,\tau}} \mathrm{hook}(a, b)$ be the hook-length product of the Young diagram $\nu_{j,\tau}$. Then the factorial-versus-hook discrepancy obeys the exact row-by-row cancellation inequality
  \begin{equation}
  \label{eq:12:young-cancel}
    \sum_{i=1}^r h(\nu_{j,\tau,i}) - \log H_{j,\tau}
    \;\le\; \mathbf{1}_{\tau \in \TT}\, h(a_\tau) + e_{j,\tau},
  \end{equation}
  where the total weight $n_{j,\tau} := \sum_i \nu_{j,\tau,i}$ splits cleanly as $n_{j,\tau} = \mathbf{1}_{\tau \in \TT} a_\tau + k w_{j,\tau} + d_{j,\tau}$ with $e_{j,\tau}$ polynomial in $N_j$ (zero super-exponential factorial excess!).
\item \emph{(Capelli branching determinant and inductive Fock norm bound.)} The integral branching map realizing the multiplicity-free $\UU(N_j) \hookrightarrow \UU(N_{j+1})$ restriction on the Fock model has explicit determinant given by the Capelli branching formula (a ratio of shifted factorials in the Gelfand--Tsetlin interlacing parameters); consequently, the branching vectors $\xi_j$ obey the inductive Fock-norm bound $\|\xi_j\|_{\mathrm{Fock}} \le C_{\mathrm{br}}^{N_j} \prod \mathrm{hook}^{1/2}$, and the terminal branching operator at layer $m$ has operator norm bounded by $\exp(O(N_m \log N_m))$.
\item \emph{(Sharp Gaussian kernel growth and single-block jet substitution.)} At every embedding $\tau \in \Phi$ and every level prime $T$, the step theta kernel satisfies the sharp Gaussian upper bound
  \begin{equation}
  \label{eq:12:sharp-gauss}
    \bigl|\theta_{\varphi_j}(g, h)\bigr|_\tau
    \;\le\; C_{\mathrm{ker}}\, \exp\!\bigl(c_1 N_j \log(k b) + k w_{j,\tau} \log\Omega_\tau\bigr)
  \end{equation}
  uniformly across the compact definite quotient, and each single-block jet operator satisfies the Fock substitution bound matching \eqref{eq:12:young-cancel}, yielding the complete analytic height-chain bound along all $m$ steps of the tower (with explicit constant $\ge 28$ absorbed into $C_0 = 3$ and $x^\varepsilon$).
\end{enumerate}
\end{theorem}
\leanref{src/Formalism/QuotientMetricDivisors/HermitianTwistDivisors/YoungDiagramTopExact.lean:212}{Formalism.DivisionValues.young\_diagram\_origin\_off}
\leanref{src/Formalism/BranchSandwichDivisors/BranchIsometryProofs/BranchMap.lean:2617}{Formalism.BranchingDeterminantCap.branch\_map\_formula}
\leanref{src/Formalism/BranchSandwichDivisors/BranchIsometryProofs/XiNorm.lean:945}{Formalism.BranchingDeterminantCap.xi\_norm\_bound}
\leanref{src/Development/AbelianVarietyTheory/Proofs/DivisorProofs/LastMapNorm.lean:395}{GlobalTheory.BranchNorm.lastMapNorm}
\leanref{src/Development/AbelianVarietyTheory/Proofs/DecoratedLatticeProofs/ChainKernelBoundAtEmbeddingAnyLevelPrime.lean:17}{GlobalTheory.Core.Ht2.KernelS.kernel\_le\_sharp}
\leanref{src/Development/AbelianVarietyTheory/Proofs/DivisorProofs/FockSubstLe.lean:119}{GlobalTheory.JetBlock.fock\_subst\_le}
\leanref{src/Formalism/LocalToGlobal/Divisibility/ConstantGeTwentyEight.lean:123}{Formalism.L2.DIV.height\_chain\_atR}

\begin{proof}
(i) For any partition $\nu = (\nu_1 \ge \cdots \ge \nu_r \ge 0)$, the hook length of box $(i, j)$ in row $i$ is $\mathrm{hook}(i, j) = (\nu_i - j) + (\nu'_j - i) + 1 \ge \nu_i - j + 1$. Taking the product over $1 \le j \le \nu_i$ in row $i$ immediately gives $\prod_{j=1}^{\nu_i} \mathrm{hook}(i, j) \ge \nu_i!$, so $\log H(\nu) \ge \sum_{i=1}^r h(\nu_i)$ for every pure partition! When a top-index twist shifts the first row by $a_\tau$ (at places $\tau \in \TT$) and introduces bounded column overlaps $d_{j,\tau}$, applying the row-by-row inequality to the sub-diagram below row~$1$ and separating the first-row excess $h(a_\tau)$ yields \eqref{eq:12:young-cancel} with exact remainder $e_{j,\tau}$ depending only on the bounded overlap height.

(ii) In the Bargmann--Fock model of the Weil representation of $\UU(N_j) \times \UU(N_{j+1})$ (\litref{Howe (1989)~\cite{Howe1989}; Fulton--Harris (1991)~\cite{FultonHarris1991}}), the polynomial inner product $\langle P, Q \rangle_{\mathrm{Fock}} = (P(\partial)\overline{Q})(0)$ satisfies $\langle z^\alpha, z^\alpha \rangle_{\mathrm{Fock}} = \alpha!$. By the Capelli identity for the determinant of the polarization matrix $\det(\sum_a z_{ia}\partial_{ja} + (r - i)\delta_{ij})$, the branching map intertwining the integral Weyl highest-weight vectors scales the Fock norm squared by the exact ratio of shifted factorials matching the hook-length product $H_{j,\tau}$! By induction on the number of branching steps (applying Cauchy--Schwarz along the Gelfand--Tsetlin tree), the norm of the branching vector $\xi_j$ cancels the factorial growth $\prod_i \nu_{j,\tau,i}!$ against $H_{j,\tau}$ via (i), bounding the terminal branching map norm cleanly.

(iii) Evaluating the definite theta series $\theta_{\varphi_j}(g, h) = \sum_{x \in \Lambda_j} P_j(x)\exp(-\pi \|x\|_\tau^2)$ via the Poisson-resummed Gaussian tail bound (using the shortest-vector lower bound on the self-dual lattice $\mathfrak{a}_0$ supplied by \cref{lem:12:rd-and-diff}(i), $\rd_E \ge 2\sqrt{d}$) bounds the sum over all nonzero lattice vectors by $O(1)$ times the leading spherical peak, while the single-block Fock substitution bound controls the derivative jets at CM points by the exact same hook-cancelled factorials, yielding \eqref{eq:12:sharp-gauss} and the global height-chain bound across all $0 \le j \le m$.
\end{proof}

\section{Four-Block Consolidation of the Eleven Arithmetic Size Inputs}
\label{sec:12:four-blocks}

In the formal archive, eleven separate arithmetic size inputs—denoted \texttt{(S1)}--\texttt{(S8)}, \texttt{(W1)}, \texttt{(W2)}, and \texttt{(G)}—together with the Siegel mass ratio, stage-zero class-number formula, and codimension-$1$ volume ratio fed into the covolume descent (\cref{ch:14}) and the definite-side assembly (\cref{ch:17}). Mirroring our four-class loss consolidation in \cref{ch:10}, we now consolidate all eleven arithmetic size inputs and mass/class-number formulas into \emph{four disjoint structural blocks}.

\newsimplification{Four-Block Structural Consolidation of All Eleven Arithmetic Size Inputs}

\begin{theorem}[Four-block arithmetic size inputs, Chowla--Selberg period cancellation, mass ratio, and stage-zero class number]
\label{thm:12:size-inputs}
Fix $0 < \varepsilon < 1/52$, an admissible host CM datum at scale $x \in [Y, Y^{3/2}]$ ($Y \ge Y_0(\psi, \varepsilon)$), and auxiliary primes $g_0, g_1, T > (k b)^8$ satisfying $2 f k < \min(g_0, g_1) - 1$. All eleven arithmetic size inputs of the period tower hold unconditionally inside the following four structural blocks:
\begin{enumerate}[label=\textup{\textbf{Block~\Roman*.}}]
\item \emph{(Archimedean chain norms, hook cancellation, and Chowla--Selberg mean periods.)}
  Let $\Omega_\tau$ ($\tau \in \Phi$) be the CM elliptic periods attached to $\mathcal{E}$ at the embeddings of $E$, and let $\log\Omega_E := \frac{1}{|\Phi|}\sum_{\tau \in \Phi} \log\Omega_\tau$ be their logarithmic mean. By the Chowla--Selberg formula (\litref{Chowla--Selberg (1967)~\cite{ChowlaSelberg1967}; Colmez (1993)~\cite{Colmez1993}}),
  \begin{equation}
  \label{eq:12:chowla-selberg}
    \log\Omega_E \;=\; -\frac{1}{2}\log(2\pi) - \frac{1}{4}\log d + \frac{1}{2}\frac{L'}{L}(0, \chi_{-d})
    \;=\; -\frac{1}{2}\log(2\pi) + \frac{1}{4}\log d - \frac{1}{2}\frac{L'}{L}(1, \chi_{-d}) - \frac{\gamma + \log(2\pi)}{2}.
  \end{equation}
  Combining \eqref{eq:12:chowla-selberg} with \eqref{eq:12:young-cancel} bounds the total chain height $h_\Pi \le h_A$ and layer averages $\langle \Sigma_j \rangle \le h_{B, j}$, and aligns the dominant period weight $k w_{j\tau}\log\Omega_E$ with the flat-determinant weight $p_j w_{j\tau}\log\Omega_E$ of \cref{ch:14} so that they cancel identically up to $|p_j - k| \le m$!

\item \emph{(Local factors, Rallis eigenvalue floor, and level-passage indices.)}
  By the convergent Rallis inner-product formula (\cref{prop:12:rallis-conj}(i)) and the formal-degree compact-dual comparison, the unramified doubling Euler product at every layer $0 \le j \le m$ obeys the explicit eigenvalue floor $l_{6, j} \ge h_{X, j}$ with total bad-prime Euler loss bounded by $-8 f \varrho^2 L_p$. Moreover, because $\min(g_0, g_1) \ge (k b)^8$ forces $3\,\mathrm{de}\,Q_0 \le 1$, the level-passage index satisfies
  \[
    L_{\mathrm{pass}} \;\le\; 1 + C_0\, m\, \lambda_0.
  \]

\item \emph{(Sharp $T$-adic valuation, twisted/glued slopes, and eight-term budget absorption.)}
  At the level prime $T$, the sharp $T$-adic $U$-operator commutation law and three-case admissibility ($p_w = T$, $p_w \in \{g_0, g_1\}$, and off-twist places) bound the glued lower-level slope directly against the covolume upper bound of \cref{ch:14} and the error budget of \cref{ch:10,ch:18}:
  \begin{equation}
  \label{eq:12:indef-boundsmu}
    \mu(R R^\vee) \;\le\; \overline{\log}(\kappa) + \sum_{j=1}^m \ell_3(j) + 2\sum_{j=1}^m\sum_{y \in \mathrm{Tw}} l_w(y) e_y F_0(y, j)
    + L_G + \sum_{j=1}^m \bigl(E_1(j) + E_2(j)\bigr) + E_3,
  \end{equation}
  where notably $E_3 = \sum 2\mathrm{wt}(y)/\ell_w(y)^2$ carries \emph{no} ramification factor $\mathrm{ram}(y)$, and the resulting eight-term budget absorbs cleanly into $(50 C_0 + 10) U_1 + 60 c_e U_2$ under $0 < \varepsilon < 1/52$.

\item \emph{(Siegel mass ratio, stage-zero class-number formula, and codimension-$1$ volume ratio.)}
  For every definite Hermitian $\OO_E$-lattice $\Lambda$ of rank $k$ over $E/E^+$, let $\mathrm{mass}_I(E, k)$ be its Siegel--Braun--Gan--Hanke--Yu mass (\litref{Siegel (1935)~\cite{Siegel1935}; Braun (1941)~\cite{Braun1941}; Gan--Hanke--Yu (2001)~\cite{GanHankeYu2001}}). By lattice-ball rank induction on binary Hermitian forms (cancelling the archimedean volume constant $c_1$ identically!), the Siegel mass ratio $\mathrm{MR}(K, M)$ is explicitly controlled across ranks $k \to k + 1$. Furthermore, since $2 f k < \min(g_0, g_1) - 1$, the supercuspidal level subgroup $\Pi(g_0, g_1)$ is neat (torsion-free!) and the stage-zero class set $\mathcal{C}_0(u)$ has exact cardinality
  \begin{equation}
  \label{eq:12:class-num}
    |\mathcal{C}_0(u)| \;=\; \mathrm{mass}_I(E, k)\,\mathrm{idx}_\Pi(E, g_0, g_1)\,(T^k - 1)^f,
  \end{equation}
  with stage-zero automorphic dimension $d(u) = d^*(u)(T^k - 1)^f$ and matching codimension-$1$ volume ratio feeding directly into the rotated wave-grid bound (\cref{ch:15}) and definite-side assembly (\cref{ch:17}).
\end{enumerate}
\end{theorem}
\leanref{src/Development/AbelianVarietyTheory/Proofs/DivisorProofs/MeanLogPeriodChowlaSelberg.lean:58}{GlobalTheory.MeanPeriodClosed.meanPeriodAt}
\leanref{src/Formalism/SharpRestriction/GlobalSchurBounds/StandardLattice.lean:281}{Formalism.MainNotions.StandardLattice.SiegelMassRatioCM}

\begin{proof}
\emph{Block~I:} Averaging the Chowla--Selberg formula~\cite{ChowlaSelberg1967,Colmez1993} over the Galois orbit $\Phi$ of embeddings of $E$ extending $\varphi_0 : K \hookrightarrow \CC$ eliminates all individual ideal-class Gamma factors and yields \eqref{eq:12:chowla-selberg} identically; note how the $-\frac{1}{2}\frac{L'}{L}(1, \chi_{-d})$ term in \eqref{eq:12:chowla-selberg} matches the $+\ell(\chi_C)$ term in $\MAIN(\ell, a, w)$ (\eqref{eq:03:main-def}) once multiplied by the rank difference! Combining \eqref{eq:12:chowla-selberg} with \eqref{eq:12:young-cancel} and \eqref{eq:12:sharp-gauss} yields $h_\Pi \le h_A$ and $\langle \Sigma_j \rangle \le h_{B, j}$.

\emph{Block~II:} In \eqref{eq:12:rallis-ip}, the unramified doubling Euler product $\prod_{v \notin S} L_v(s_0 + 1/2, \pi_v \times \chi_v) / b_v(s_0, \chi_v)$ is bounded below by comparing against $\zeta_E(2)^{-1}$ and extracting the finite bad-prime loss $-8f\varrho^2 L_p$. For the level-passage index, the ratio of compact-open subgroup indices at $g_0, g_1$ is bounded by $\prod (1 + O(g_i^{-1}))^{m\lambda_0}$; since $\min(g_0, g_1) \ge (k b)^8$, $3\,\mathrm{de}\,Q_0 \le 1$, yielding $L_{\mathrm{pass}} \le 1 + C_0 m \lambda_0$.

\emph{Block~III:} Commuting the $T$-adic Hecke operator $U_T$ with the level-$T$ evaluation yields the sharp $T$-adic valuation shift matching \cref{prop:10:level-prime}(i); summing over conjugate pairs of twisted places $(y, \bar{y})$ produces the factor $2$ in front of $\sum_{y \in \mathrm{Tw}} l_w(y) e_y F_0(y, j)$ in \eqref{eq:12:indef-boundsmu}, while the error term $E_3$ depends solely on $\mathrm{wt}(y)/\ell_w(y)^2$ without any $\mathrm{ram}(y)$ multiplier. Applying \cref{thm:10:four-class-losses} and \cref{thm:18:twisted-integrality} absorbs all eight error summands into $(50C_0 + 10)U_1 + 60 c_e U_2$.

\emph{Block~IV:} The Siegel--Braun--Gan--Hanke--Yu mass formula~\cite{Siegel1935,Braun1941,GanHankeYu2001} expresses $\mathrm{mass}_I(E, k)$ as a product of special values $L(r, \omega_{E/E^+}^r)$ ($1 \le r \le k$) times local lattice densities $\alpha_v(\Lambda_v)^{-1}$. Passing from rank $k-1$ to $k$ by binary Hermitian lattice-ball counting cancels the universal archimedean constant $c_1$ identically. Finally, any nontrivial root of unity $\zeta \in \mu(E)$ that lies in the principal congruence subgroup at $g_i$ must satisfy $g_i \mid \Nm(\zeta - 1)$, which is impossible once $2 f k < g_i - 1$ (since $w_E \le 8 f^2$ and every nontrivial $\zeta$ has minimal polynomial of bounded degree relative to $g_i$, and in fact the order of any torsion element in $\UU_k(\OO_E)$ has Euler totient $\le 2 f k < g_i - 1$, forcing every stabilizer in the level subgroup $\Pi(g_0, g_1)$ to be trivial!). Hence every class in $\mathcal{C}_0(u)$ has trivial stabilizer, so the weighted mass equals the unweighted cardinality \eqref{eq:12:class-num}, with the factor $(T^k - 1)^f$ coming from the $f$ split places above $T$ in $E^+$.
\end{proof}
